Este capítulo aplica ao subsistema descrito no \refCap{3_Implementacao} os critérios de verificação apresentados na seção~\ref{sec:2.1}. A exposição se divide em duas partes, a primeira responde se o projeto faz o que a especificação determina, e tem por instrumento a bateria de testes automatizados. Por outro lado, a segunda responde quanto o subsistema custa e quanto economiza, e tem por instrumento os contadores de desempenho. Em síntese, as duas se apoiam no mesmo ambiente de execução, mas medem coisas distintas já que uma produz um veredito binário por caso e a outra produz números que só significam algo quando acompanhados da descrição da carga que os gerou.

\section{Ambiente de verificação e resultados gerais}\label{sec:4.1}

Os testes são escritos com ScalaTest e ChiselTest, sobre Scala 2.13.14, Chisel 3.6.1 e chiseltest 0.6.1, com sbt 1.12.8. O simulador empregado é o Treadle, interpretador escrito em Scala que acompanha a série 3.x do Chisel e executa o circuito elaborado dentro da própria máquina virtual. O Verilator é alternativa disponível, mas não é o simulador padrão do projeto: ele aparece em um único caso, marcado como lento na suíte de política de reposição, e é ativado por variável de ambiente. Como a integração contínua não instala o Verilator nem define essa variável, todos os resultados deste capítulo provêm do Treadle.

Por conseguinte, a integração contínua executa, em tarefas separadas, a verificação de formatação e a bateria completa de testes a cada alteração. A execução de referência, de 23 de novembro de 2026, compreendeu trinta suítes e cento e oitenta casos: cento e setenta e nove aprovados, nenhum reprovado e um cancelado, em dois minutos e vinte e quatro segundos. O caso cancelado é o que executa um programa em hexadecimal e exige que o caminho do arquivo seja fornecido por variável de ambiente.

\section{Organização dos testes e cobertura das operações}\label{sec:4.2}

As trinta suítes se distribuem em três grupos: setenta e seis casos sobre o subsistema de cache, cinquenta e quatro sobre a integração com o pipeline e cinquenta sobre os demais módulos do núcleo. A \refTab{tab:4.1} reúne as suítes de cache, que são as que interessam a este trabalho.

{\small
\begin{longtable}{@{}>{\raggedright\arraybackslash}p{(\textwidth - 4\tabcolsep) * \real{0.2729}}>{\raggedright\arraybackslash}p{(\textwidth - 4\tabcolsep) * \real{0.0915}}>{\raggedright\arraybackslash}p{(\textwidth - 4\tabcolsep) * \real{0.6356}}@{}}
\caption[Suítes de teste do subsistema de cache e respectiva cobertura]{Suítes de teste do subsistema de cache e respectiva cobertura (Fonte: autoria própria)}\label{tab:4.1}\\
\hline
\textbf{Suíte} & \textbf{Casos} & \textbf{Cobertura} \\
\hline
\endfirsthead
\hline
\textbf{Suíte} & \textbf{Casos} & \textbf{Cobertura} \\
\hline
\endhead
\hline
\endlastfoot
CacheTest & 8 & geometria derivada, soma dos campos do endereço, independência entre as duas caches \\
CacheArraysTest & 4 & estado inicial, escrita e leitura de linha, máscara de palavra, habilitação de metadados \\
CacheDirtyBitTest & 3 & existência condicional do bit de sujeira, leitura constante em zero na cache de instruções \\
CacheFsmTest & 8 & máquina de estados: repouso, falta de leitura, falta de escrita, despejo, drenagem, acesso de byte e meia-palavra \\
CacheHitTest & 5 & caminho de acerto nas quatro vias, alternância entre vias, escrita por byte e meia-palavra \\
CacheLruTest & 4 & ordem de reposição, despejo por via, independência por conjunto, modelo de referência \\
CacheMainMemoryTest & 9 & pré-carga, visibilidade entre as duas caches, arbitragem, limite de capacidade, reset \\
CacheMemoryLatencyTest & 7 & latência 1, 3 e 6; arbitragem e aborto; falta simultânea; consistência \\
CacheFaultTest & 5 & canal de erro no preenchimento, no despejo e na drenagem \\
CachePerfTest & 2 & os seis contadores, na cache isolada e no coordenador \\
CacheGeometryTest & 17 & validação de configurações e varredura de treze geometrias contra modelo \\
CacheSweepTraceTest & 4 & determinismo do trace e independência em relação à geometria \\
\end{longtable}
}

Diante desse panorama, as cinco operações que definem o comportamento funcional de uma cache de dados têm verificação direta, e cada uma é confirmada por mais de uma suíte, em níveis distintos de abstração: a máquina de estados isolada, o caminho de acerto em todas as vias, a memória com latência e a varredura contra modelo. A \refTab{tab:4.2} reúne essa correspondência.

{\small
\begin{longtable}{@{}>{\raggedright\arraybackslash}p{(\textwidth - 2\tabcolsep) * \real{0.2500}}>{\raggedright\arraybackslash}p{(\textwidth - 2\tabcolsep) * \real{0.7500}}@{}}
\caption[Cobertura das operações de cache e o que cada conjunto de testes comprova]{Cobertura das operações de cache e o que cada conjunto de testes comprova (Fonte: autoria própria)}\label{tab:4.2}\\
\hline
\textbf{Operação} & \textbf{O que os testes comprovam} \\
\hline
\endfirsthead
\hline
\textbf{Operação} & \textbf{O que os testes comprovam} \\
\hline
\endhead
\hline
\endlastfoot
Leitura com acerto & custa dois ciclos e não acessa a memória; acerta nas quatro vias e em todas as palavras da linha; a cache de instruções devolve a palavra inteira \\
Leitura com falta & o preenchimento lê exatamente uma linha e nenhuma escrita é emitida; o custo cresce com a latência; faltas simultâneas das duas caches são serializadas \\
Escrita com acerto & grava apenas a palavra alvo e marca sujeira somente na linha escrita; acesso de byte e de meia-palavra faz mesclagem; a memória não é atualizada \\
Escrita com falta & a linha é trazida, a escrita é aplicada na cache e a linha fica suja, sem tráfego de escrita à memória \\
Substituição de linha suja & o despejo gera exatamente um despejo com o conteúdo correto; linha limpa não gera tráfego; a vítima segue a ordem de recência por conjunto \\
\end{longtable}
}

Três verificações dessa tabela merecem comentário, porque sustentam afirmações feitas em capítulos anteriores. A primeira é que a igualdade entre acertos e faltas, de um lado, e o número de pedidos, de outro, é verificada explicitamente, o que materializa a convenção fixada na seção~\ref{sec:3.10} de que o acerto que atende o pedido reavaliado após um preenchimento não é contado como acerto. A segunda é que o custo de cinco faltas, sendo quatro limpas e uma com despejo, é conferido como onze ciclos de falta, o que confirma a fórmula da seção~\ref{sec:3.8} no caso em que a latência é zero. A terceira é que, com latência três, a cache de instruções que solicita uma linha em processo de despejo só é atendida depois que o despejo termina, o que verifica a regra de prioridade da cache de dados fixada na seção~\ref{sec:3.6}.

Uma lacuna de cobertura se dá no contexto da política de reposição aleatória, pois essa é exercitada apenas na varredura, onde os dados são conferidos contra o modelo, mas os números de acerto, falta e despejo não têm modelo de referência correspondente, porque o modelo implementa LRU e FIFO e a sequência do registrador de deslocamento é reproduzível mas não trivial de antecipar. Em consequência, o resultado da política aleatória na seção~\ref{sec:4.5} verifica que o dado lido está correto, e não que a escolha da vítima obedeceu à política.

\section{A verificação do mecanismo de consistência}\label{sec:4.3}

O mecanismo de consistência recebe tratamento próprio porque é o único ponto do subsistema em que a correção depende de uma ordem entre operações, e não apenas do resultado de cada uma. A \refTab{tab:4.3} reúne os testes que o exercitam.

{\small
\begin{longtable}{@{}>{\raggedright\arraybackslash}p{(\textwidth - 4\tabcolsep) * \real{0.2222}}>{\raggedright\arraybackslash}p{(\textwidth - 4\tabcolsep) * \real{0.3333}}>{\raggedright\arraybackslash}p{(\textwidth - 4\tabcolsep) * \real{0.4444}}@{}}
\caption[Testes do mecanismo de consistência e o que cada um comprova]{Testes do mecanismo de consistência e o que cada um comprova (Fonte: autoria própria)}\label{tab:4.3}\\
\hline
\textbf{Teste} & \textbf{Cenário} & \textbf{O que comprova} \\
\hline
\endfirsthead
\hline
\textbf{Teste} & \textbf{Cenário} & \textbf{O que comprova} \\
\hline
\endhead
\hline
\endlastfoot
Integração, código automodificável & sub-rotina reescrita por armazenamento após a primeira chamada, seguida da instrução de sincronização e de nova chamada & a segunda chamada devolve o valor novo, com latência zero e com latência quatro \\
Latência de memória & dois armazenamentos com alocação, sincronização, busca pela cache de instruções & a cache de instruções lê os valores gravados pela cache de dados; a palavra não escrita permanece zero \\
Tráfego aleatório & sequência de leituras e escritas, depois sincronização & todo o conteúdo escrito está na memória e é visível à cache de instruções \\
Contadores & armazenamento, sincronização, leitura dos contadores & os despejos de drenagem são contados separadamente dos despejos por falta \\
Máquina de estados & uma linha suja entre linhas limpas & apenas a linha suja gera tráfego; ela permanece residente, mas limpa \\
Varredura & drenagem após o trace, nas treze configurações & a drenagem grava exatamente as linhas sujas, e nenhum dado se perde \\
Canal de erro & memória recusa o despejo da drenagem & a operação aborta sem travar e assinala a falha; nova drenagem sem erro conclui e limpa o sinal \\
\end{longtable}
}

Um ponto a se levantar é que, no caso de código automodificável, a sub-rotina é chamada, devolve um valor, é reescrita por um armazenamento, que deixa a linha suja na cache de dados, a instrução de sincronização é executada e a sub-rotina é chamada de novo. Se a drenagem não ocorresse, a memória ainda conteria a versão antiga e, caso a invalidação não ocorresse, a cache de instruções ainda conteria a versão antiga. Como o teste exige o valor novo nas duas latências, ele falharia em qualquer das duas omissões, e é essa conjunção que verifica a obrigação da seção~\ref{sec:3.9} como um todo, e não cada operação isoladamente.

Todavia, uma lacuna observada é a sinalização de falha da sincronização, a qual é verificada apenas na cache isolada, não no coordenador nem no pipeline, o que decorre de a saída de falha do coordenador apenas repassar o sinal da cache de dados, sem lógica própria a verificar. O custo em ciclos da instrução foi objeto de instrumentação própria e é apresentado nas seções~\ref{sec:4.3.1} a~\ref{sec:4.3.3}.

\subsection{Custo na cache isolada}\label{sec:4.3.1}

A medição na cache isolada liga o coordenador a uma memória de latência configurável e executa três passos. Primeiro, sujam-se N linhas por meio de armazenamentos em endereços de linhas distintas, de modo que cada um deixe uma linha suja sem provocar despejo. Em seguida, ergue-se o pedido de sincronização e contam-se os ciclos de relógio até o sinal de conclusão. Por fim, confere-se o resultado tanto pela expressão de custo quanto pelo contador de despejos de drenagem, que deve registrar exatamente as N linhas. Como a contagem é feita em ciclos de relógio sobre um modelo determinístico, a medida não tem ruído, e o resultado é uma verificação da máquina de estados.

A \refTab{tab:4.4} apresenta os resultados na configuração de 1 KiB, com dezesseis conjuntos, e a \refTab{tab:4.5} os da configuração de base, com cento e vinte e oito conjuntos. Todas as medições desta seção foram feitas no modo de arranjo com leitura combinacional, sem blocos de memória dedicados, e os números não valem para a configuração sintetizada, em que o custo por conjunto e por linha suja é maior, conforme a seção~\ref{sec:3.9}.

{\small
\begin{longtable}{@{}>{\raggedright\arraybackslash}p{(\textwidth - 8\tabcolsep) * \real{0.2100}}>{\raggedright\arraybackslash}p{(\textwidth - 8\tabcolsep) * \real{0.2100}}>{\raggedright\arraybackslash}p{(\textwidth - 8\tabcolsep) * \real{0.2300}}>{\raggedright\arraybackslash}p{(\textwidth - 8\tabcolsep) * \real{0.1700}}>{\raggedright\arraybackslash}p{(\textwidth - 8\tabcolsep) * \real{0.1800}}@{}}
\caption[Custo da sincronização na cache isolada de 1 KiB, em ciclos, por latência da memória e número de linhas sujas]{Custo da sincronização na cache isolada de 1 KiB, em ciclos, por latência da memória e número de linhas sujas (Fonte: autoria própria, a partir da execução de 23 de novembro de 2026)}\label{tab:4.4}\\
\hline
\textbf{Linhas sujas} & \textbf{L = 0} & \textbf{L = 1} & \textbf{L = 3} & \textbf{L = 6} \\
\hline
\endfirsthead
\hline
\textbf{Linhas sujas} & \textbf{L = 0} & \textbf{L = 1} & \textbf{L = 3} & \textbf{L = 6} \\
\hline
\endhead
\hline
\endlastfoot
0 & 20 & 20 & 20 & 20 \\
1 & 22 & 23 & 25 & 28 \\
2 & 24 & 26 & 30 & 36 \\
64 & 148 & 212 & 340 & 532 \\
\end{longtable}
}

A \refFig{fig:4.1} representa os dados da \refTab{tab:4.4} e mostra que as quatro curvas partem do mesmo patamar quando não há linha suja, separando-se à medida que o número de linhas a gravar cresce.

\figuraFonte[!ht]{figura_4_1}{Custo da instrução de sincronização em ciclos, por latência da memória e número de linhas sujas}{autoria própria, a partir da execução de 23 de novembro de 2026}{fig:4.1}{width=\textwidth}

{\small
\begin{longtable}{@{}>{\raggedright\arraybackslash}p{(\textwidth - 4\tabcolsep) * \real{0.4689}}>{\raggedright\arraybackslash}p{(\textwidth - 4\tabcolsep) * \real{0.2913}}>{\raggedright\arraybackslash}p{(\textwidth - 4\tabcolsep) * \real{0.2398}}@{}}
\caption[Custo da sincronização na configuração de base, com 128 conjuntos e latência três, medido e previsto]{Custo da sincronização na configuração de base, com 128 conjuntos e latência três, medido e previsto (Fonte: autoria própria, a partir da execução de 23 de novembro de 2026)}\label{tab:4.5}\\
\hline
\textbf{Linhas sujas} & \textbf{Medido} & \textbf{Previsto} \\
\hline
\endfirsthead
\hline
\textbf{Linhas sujas} & \textbf{Medido} & \textbf{Previsto} \\
\hline
\endhead
\hline
\endlastfoot
0 & 132 & 132 \\
1 & 137 & 137 \\
2 & 142 & 142 \\
512 & 2.692 & 2.692 \\
\end{longtable}
}

Em todas as vinte combinações ensaiadas o valor medido coincidiu exatamente com o previsto, sem divergência de um único ciclo. A expressão que os dados confirmam é ciclos = numSets + 4 + $\Sigma$ (L + 2), em que cada linha suja responde por um ciclo de varredura e por L + 1 ciclos de gravação. Como consequência, a expressão da seção~\ref{sec:4.8} fica verificada, isto é, o circuito e o cálculo são descrições independentes que convergem em vinte pontos. Além disso, os dados mostram quatro ciclos a mais, correspondentes ao handshake entre a máquina de estados do coordenador e a da cache. O mínimo, portanto, não é 128 ciclos na configuração de base, e sim 132.

O pior caso construído da configuração de base, com todas as 512 linhas sujas e latência três, custa 2.692 ciclos, dos quais 2.560 correspondem às gravações e 132 à varredura. Com poucas linhas sujas, o custo é dominado pela varredura, em 1 KiB, uma única linha suja acrescenta dois ciclos a um piso de vinte.

\subsection{Custo no pipeline}\label{sec:4.3.2}

A medição no pipeline executa um programa que lê o contador de ciclos, executa a instrução de sincronização e lê o contador novamente. Como a janela medida inclui as próprias instruções de leitura, um segundo programa idêntico, com uma instrução vazia no lugar da sincronização, calibra esse custo e é subtraído.

A \refTab{tab:4.6} apresenta os resultados, com o custo no pipeline, o custo correspondente na cache isolada e a diferença entre os dois. Como na seção anterior, as medições foram feitas no modo de arranjo sem blocos de memória dedicados, e os valores não se transferem à configuração sintetizada. A medição grava os resultados em dois arquivos do repositório, um para a cache isolada e outro para o pipeline, de modo que os números das Tabelas 4.4, 4.5 e 4.6 são reprodutíveis.

{\small
\begin{longtable}{@{}>{\raggedright\arraybackslash}p{(\textwidth - 8\tabcolsep) * \real{0.2100}}>{\raggedright\arraybackslash}p{(\textwidth - 8\tabcolsep) * \real{0.2100}}>{\raggedright\arraybackslash}p{(\textwidth - 8\tabcolsep) * \real{0.2300}}>{\raggedright\arraybackslash}p{(\textwidth - 8\tabcolsep) * \real{0.1700}}>{\raggedright\arraybackslash}p{(\textwidth - 8\tabcolsep) * \real{0.1800}}@{}}
\caption[Custo da sincronização no pipeline, em ciclos, por latência e número de linhas sujas]{Custo da sincronização no pipeline, em ciclos, por latência e número de linhas sujas (Fonte: autoria própria, a partir da execução de 23 de novembro de 2026)}\label{tab:4.6}\\
\hline
\textbf{Latência} & \textbf{Linhas sujas} & \textbf{Custo no pipeline} & \textbf{Custo na cache} & \textbf{Overhead} \\
\hline
\endfirsthead
\hline
\textbf{Latência} & \textbf{Linhas sujas} & \textbf{Custo no pipeline} & \textbf{Custo na cache} & \textbf{Overhead} \\
\hline
\endhead
\hline
\endlastfoot
0 & 0 & 26 & 20 & 6 \\
0 & 1 & 28 & 22 & 6 \\
0 & 2 & 30 & 24 & 6 \\
0 & 64 & 154 & 148 & 6 \\
1 & 0 & 27 & 20 & 7 \\
1 & 1 & 30 & 23 & 7 \\
1 & 2 & 33 & 26 & 7 \\
1 & 64 & 219 & 212 & 7 \\
3 & 0 & 29 & 20 & 9 \\
3 & 1 & 34 & 25 & 9 \\
3 & 2 & 39 & 30 & 9 \\
3 & 64 & 349 & 340 & 9 \\
6 & 0 & 32 & 20 & 12 \\
6 & 1 & 40 & 28 & 12 \\
6 & 2 & 48 & 36 & 12 \\
6 & 64 & 544 & 532 & 12 \\
\end{longtable}
}

A diferença entre o custo no pipeline e o custo na cache isolada é constante em L + 6 ciclos e não depende do número de linhas sujas. Essa constância admite uma interpretação que se apoia na \refTab{tab:3.10}. Depois da invalidação, a cache de instruções está integralmente vazia, e a primeira busca que incide sobre cada bloco é uma falta compulsória, cujo custo é L + 4 ciclos. Somados dois ciclos de encerramento das duas máquinas, chega-se a L + 6. Portanto, o overhead do pipeline equivale ao custo de uma falta limpa mais dois ciclos. Trata-se de interpretação consistente com todas as medições.

O custo marginal por linha suja, isolado pela diferença entre execuções, é de L + 2, isto é, o mesmo valor obtido na cache isolada. A convergência entre os dois níveis de medição reforça a interpretação de que o custo do pipeline é a soma do custo da drenagem com um termo fixo de retomada.

\subsection{Consequências para a geometria}\label{sec:4.3.3}

Um resultado da medição não se refere à sincronização em si, mas à parametrização. A varredura percorre todos os conjuntos, e não apenas os que contêm linhas sujas. Em consequência, uma cache maior paga mais por sincronização mesmo quando não há nada a gravar, a saber, vinte ciclos na configuração de 1 KiB contra 132 na de 8 KiB, com zero linhas sujas, o que corresponde a um acréscimo de 112 ciclos.

Diante do pior caso de 2.692 ciclos, 112 ciclos adicionais representam cerca de 4\% do total. Diante do caso em que a cache está limpa e a sincronização apenas varre, representam mais de seis vezes o custo. Como a instrução é rara e concentrada em carregamento de código e em código automodificável, o impacto agregado é pequeno, mas o resultado estabelece que a geometria adotada, que a seção~\ref{sec:4.6} justificou por saturação da taxa de acerto, tem um custo próprio na operação de consistência, e esse custo cresce com a capacidade.

\section{Metodologia de medição}\label{sec:4.4}

Cada cache dispõe de seis contadores de sessenta e quatro bits, agrupados em um agregado e expostos pelo coordenador ao pipeline. Uma entrada de limpeza zera todos simultaneamente, o que permite delimitar a janela de medição sem reiniciar a simulação. Os contadores são espelhados em registradores de controle e status numerados de três a doze na faixa de contadores de desempenho, em modo somente leitura, com escritas ignoradas, conforme a especificação da arquitetura para esses registradores, ao lado dos contadores de ciclos e de instruções aposentadas. A \refTab{tab:4.7} descreve cada um.

{\small
\begin{longtable}{@{}>{\raggedright\arraybackslash}p{(\textwidth - 4\tabcolsep) * \real{0.2226}}>{\raggedright\arraybackslash}p{(\textwidth - 4\tabcolsep) * \real{0.2226}}>{\raggedright\arraybackslash}p{(\textwidth - 4\tabcolsep) * \real{0.5548}}@{}}
\caption[Contadores de desempenho, respectivo registrador e definição]{Contadores de desempenho, respectivo registrador e definição (Fonte: autoria própria)}\label{tab:4.7}\\
\hline
\textbf{Contador} & \textbf{Registrador} & \textbf{O que conta} \\
\hline
\endfirsthead
\hline
\textbf{Contador} & \textbf{Registrador} & \textbf{O que conta} \\
\hline
\endhead
\hline
\endlastfoot
acertos & 3 (instruções), 6 (dados) & pedidos atendidos por acerto, excluído o ciclo posterior ao preenchimento \\
faltas & 4 (instruções), 7 (dados) & pedidos que falharam; o pedido refeito após uma falha conta nova falta \\
ciclos de falta & 5 (instruções), 8 (dados) & ciclos gastos no tratamento da falta \\
despejos & 9 (dados) & linhas sujas gravadas por substituição \\
despejos de drenagem & 10 (dados) & linhas sujas gravadas pela sincronização, contadas à parte \\
erros & 11 (instruções), 12 (dados) & preenchimentos recusados pela memória \\
\end{longtable}
}

Nesse sentido, os ciclos até o acerto são contados pelo próprio teste, que avança o relógio e conta as transições, já os ciclos de falta vêm do contador, que a máquina de estados incrementa apenas nos estados de tratamento. São unidades medidas em pontos diferentes do circuito e não devem ser somadas sem essa ressalva.

Os contadores são verificados em três níveis: na cache isolada, nas saídas do coordenador e nos registradores de controle e status lidos pelo pipeline. Cabe notar que apenas três desses contadores são lidos pela via dos registradores no teste de integração, enquanto os demais são verificados na porta da cache ou nas saídas do coordenador.

Cada configuração da varredura de geometria, que produz os números da seção~\ref{sec:4.5}, é uma cache de dados isolada, sobre uma memória de trinta e dois kibibytes com latência zero, variando um eixo por vez a partir da configuração de base. O trace é único e fixo: semente quarenta e dois, mil e quinhentos acessos, área de endereçamento de vinte e quatro kibibytes, rajadas de uma a quatro palavras, vinte e cinco por cento de escritas, com sessenta por cento das rajadas concentradas em uma região quente de quatro kibibytes, vinte e cinco por cento em uma região morna de doze kibibytes e quinze por cento distribuídos pela área total.

O teste impõe duas condições que dão credibilidade aos números. A primeira é que a soma de acertos e faltas seja igual ao número de acessos, o que detecta qualquer erro de contagem. A segunda é que acertos, faltas e despejos coincidam com o modelo de referência para as políticas LRU e FIFO, nas treze configurações.

\section{Taxa de acerto e de falta por configuração}\label{sec:4.5}

Os resultados da varredura se distribuem em quatro eixos. As Tabelas~\ref{tab:4.8} e~\ref{tab:4.9} apresentam os dois eixos com efeito mensurável, enquanto os outros dois são comentados em seguida.

{\small
\begin{longtable}{@{}>{\raggedright\arraybackslash}p{(\textwidth - 12\tabcolsep) * \real{0.1786}}>{\raggedright\arraybackslash}p{(\textwidth - 12\tabcolsep) * \real{0.1419}}>{\raggedright\arraybackslash}p{(\textwidth - 12\tabcolsep) * \real{0.1319}}>{\raggedright\arraybackslash}p{(\textwidth - 12\tabcolsep) * \real{0.1319}}>{\raggedright\arraybackslash}p{(\textwidth - 12\tabcolsep) * \real{0.1319}}>{\raggedright\arraybackslash}p{(\textwidth - 12\tabcolsep) * \real{0.1419}}>{\raggedright\arraybackslash}p{(\textwidth - 12\tabcolsep) * \real{0.1419}}@{}}
\caption[Taxa de acerto e de falta no eixo de capacidade, com quatro vias, linha de dezesseis bytes e LRU]{Taxa de acerto e de falta no eixo de capacidade, com quatro vias, linha de dezesseis bytes e LRU (Fonte: autoria própria, a partir da execução de 23 de novembro de 2026)}\label{tab:4.8}\\
\hline
\textbf{Capacidade} & \textbf{Conjuntos} & \textbf{Acertos} & \textbf{Faltas} & \textbf{Despejos} & \textbf{Taxa de acerto} & \textbf{Taxa de falta} \\
\hline
\endfirsthead
\hline
\textbf{Capacidade} & \textbf{Conjuntos} & \textbf{Acertos} & \textbf{Faltas} & \textbf{Despejos} & \textbf{Taxa de acerto} & \textbf{Taxa de falta} \\
\hline
\endhead
\hline
\endlastfoot
1 KiB & 16 & 778 & 722 & 261 & 51,9\% & 48,1\% \\
2 KiB & 32 & 873 & 627 & 207 & 58,2\% & 41,8\% \\
4 KiB & 64 & 981 & 519 & 126 & 65,4\% & 34,6\% \\
8 KiB & 128 & 1.043 & 457 & 16 & 69,5\% & 30,5\% \\
16 KiB & 256 & 1.044 & 456 & 1 & 69,6\% & 30,4\% \\
\end{longtable}
}

A \refFig{fig:4.2} representa esses mesmos dados em forma de curva, e evidencia a saturação entre 4 KiB e 8 KiB.

\figuraFonte[!ht]{figura_4_2}{Taxa de acerto em função da capacidade da cache, com quatro vias, linha de dezesseis bytes e LRU}{autoria própria, a partir da execução de 23 de novembro de 2026}{fig:4.2}{width=0.85\textwidth}

{\small
\begin{longtable}{@{}>{\raggedright\arraybackslash}p{(\textwidth - 12\tabcolsep) * \real{0.1786}}>{\raggedright\arraybackslash}p{(\textwidth - 12\tabcolsep) * \real{0.1419}}>{\raggedright\arraybackslash}p{(\textwidth - 12\tabcolsep) * \real{0.1319}}>{\raggedright\arraybackslash}p{(\textwidth - 12\tabcolsep) * \real{0.1319}}>{\raggedright\arraybackslash}p{(\textwidth - 12\tabcolsep) * \real{0.1319}}>{\raggedright\arraybackslash}p{(\textwidth - 12\tabcolsep) * \real{0.1419}}>{\raggedright\arraybackslash}p{(\textwidth - 12\tabcolsep) * \real{0.1419}}@{}}
\caption[Taxa de acerto e de falta no eixo de tamanho de linha, com oito kibibytes, quatro vias e LRU]{Taxa de acerto e de falta no eixo de tamanho de linha, com oito kibibytes, quatro vias e LRU (Fonte: autoria própria, a partir da execução de 23 de novembro de 2026)}\label{tab:4.9}\\
\hline
\textbf{Linha} & \textbf{Conjuntos} & \textbf{Acertos} & \textbf{Faltas} & \textbf{Despejos} & \textbf{Taxa de acerto} & \textbf{Taxa de falta} \\
\hline
\endfirsthead
\hline
\textbf{Linha} & \textbf{Conjuntos} & \textbf{Acertos} & \textbf{Faltas} & \textbf{Despejos} & \textbf{Taxa de acerto} & \textbf{Taxa de falta} \\
\hline
\endhead
\hline
\endlastfoot
8 B & 256 & 799 & 701 & 6 & 53,3\% & 46,7\% \\
16 B & 128 & 1.043 & 457 & 16 & 69,5\% & 30,5\% \\
32 B & 64 & 1.173 & 327 & 41 & 78,2\% & 21,8\% \\
64 B & 32 & 1.263 & 237 & 63 & 84,2\% & 15,8\% \\
\end{longtable}
}

No eixo de associatividade, mantidos oito kibibytes, linha de dezesseis bytes e LRU, a taxa de acerto é de 67,2\% com uma via, 68,9\% com duas, 69,5\% com quatro e 69,6\% com oito. No eixo de política, mantidos oito kibibytes, quatro vias e linha de dezesseis bytes, a taxa é de 69,5\% com LRU, 69,3\% com escolha aleatória e 69,0\% com FIFO.

\section{Análise dos resultados}\label{sec:4.6}

O primeiro resultado a destacar é que os eixos de capacidade e de associatividade saturam exatamente na configuração adotada no projeto. Na capacidade, os incrementos sucessivos são de 6,3, 7,2, 4,1 e 0,1 ponto percentual, isto é, dobrar de oito para dezesseis kibibytes rende um único acerto adicional em mil e quinhentos acessos. Na associatividade, os incrementos são de 1,7, 0,6 e 0,1 ponto, em que dobrar de quatro para oito vias rende igualmente um acerto.

Esse resultado tem uma consequência para o \refCap{3_Implementacao}. A justificativa da geometria na seção~\ref{sec:3.2} se apoiava em dados de síntese de terceiros, que indicavam o custo em área de cada parâmetro, mas não no comportamento do subsistema sob carga. A varredura fornece o complemento que faltava, sendo que dentro da faixa avaliada e sob este trace, 8 KiB e 4 vias estão no ponto em que o ganho marginal já se anulou, de modo que o custo adicional de 16 KiB ou de 8 vias não se traduziria em acerto.

Há um segundo efeito, mais informativo que a taxa de acerto. O número de despejos cai de 261 para 16 entre 1 e 8 KiB, e para 1 em 16 KiB. Traduzido em tráfego, isso significa que a cache de base já captura quase todas as linhas que seriam despejadas sujas. Dessa forma, o argumento a favor da configuração adotada não é apenas que ela acerta tanto quanto uma maior, mas que ela já elimina praticamente todo o tráfego de despejo que uma cache menor produziria, e o tráfego de despejo é justamente o que a política de escrita de volta torna caro, conforme a seção~\ref{sec:2.5.3}.

O eixo de tamanho de linha produz o resultado de maior amplitude, e é o único em que a configuração adotada não é a melhor. A taxa de acerto cresce de 53,3\% com oito bytes para 84,2\% com sessenta e quatro bytes, e a linha de sessenta e quatro supera a de dezesseis em 14,7 pontos percentuais, com 220 faltas a menos. Lido apenas pela taxa de acerto, esse resultado condenaria a escolha da seção~\ref{sec:3.2}. Lido em bytes transferidos, inverte. O tráfego total, soma do que entra por preenchimento e do que sai por despejo, é de 5.656 bytes com linha de oito, 7.568 com dezesseis, 11.776 com trinta e dois e 19.200 com sessenta e quatro. A linha de sessenta e quatro bytes move 2,5 vezes mais bytes que a de dezesseis para evitar 220 faltas.

Em número de ocorrências, os despejos crescem de 6 para 63 ao longo do eixo; em bytes, crescem de 48 para 4.032. Comparar ocorrências entre configurações de linha diferente mede a granularidade da transferência, não o custo. Toda a análise de tráfego deste capítulo é feita em bytes por essa razão.

Portanto, a leitura conjunta dos dois eixos sustenta a seguinte conclusão sobre a configuração adotada. Ela não maximiza a taxa de acerto, e não poderia: a linha de sessenta e quatro bytes acerta 14,7 pontos mais. Ela minimiza o tráfego dentro da faixa de acerto aceitável, e essa é uma escolha de projeto coerente com um núcleo cujo objetivo é executar um sistema operacional sobre uma memória simulada de largura de banda fixa. Um sistema em que cada transferência custa o mesmo tempo, qualquer que seja o seu tamanho, favoreceria a linha maior, já um sistema em que a transferência é proporcional ao volume transferido favorece a menor. A seção~\ref{sec:4.7} retoma essa distinção.

Por fim, o eixo de política confirma a ordenação prevista na seção~\ref{sec:2.5.4}, com a LRU à frente da escolha aleatória, que por sua vez fica à frente da FIFO, mas com margem pequena. A diferença entre a primeira e a última é de oito acertos em mil e quinhentos acessos, ou 0,5 ponto percentual. O que os dados sustentam é que a ordenação prevista se confirma e que a magnitude é pequena sob esta carga, de modo que a previsibilidade e a testabilidade permanecem como justificativas mais robustas que o ganho de acerto. Cabe notar ainda que uma diferença de oito acertos em mil e quinhentos não é estatisticamente distinguível de ruído, o que reforça a leitura de que a política é o parâmetro de menor impacto entre os quatro.

\section{Limitações da avaliação}\label{sec:4.7}

Cinco limitações delimitam o alcance dos números apresentados.

A primeira é a natureza do trace. Ele é sintético e construído para exibir localidade espacial e temporal: sessenta por cento das rajadas em quatro kibibytes, rajadas sequenciais, vinte e cinco por cento de escritas. Os resultados valem para esse padrão, e não para um programa real. Um trace com predomínio de saltos inverteria a conclusão do eixo de linha, porque linhas maiores passariam a trazer dados que nunca seriam usados. Assim, a curva da seção~\ref{sec:4.5} é uma curva de sensibilidade a parâmetros sob uma carga declarada, e não uma previsão de desempenho.

A segunda é o objeto medido. Cada configuração da varredura é uma cache de dados isolada, sobre memória de latência zero. O primeiro nível completo tem duas caches, com arbitragem entre elas, e o custo de falta depende da latência, que aqui é nula. Dessa maneira, os números medem o efeito da geometria sobre a taxa de acerto, e não o custo em ciclos do subsistema em operação em um sistema real.

A terceira é o tamanho da amostra. Mil e quinhentos acessos são suficientes para curvas com diferenças de vários pontos percentuais, como as dos eixos de capacidade e de linha, e insuficientes para os eixos em que a diferença é de poucos acertos, como o de política e o trecho saturado do de associatividade. Nesses casos o resultado indica ordem de grandeza, não distinção.

A quarta é o escopo das medições. Não há comparação entre execução com e sem cache para um mesmo programa, porque o pipeline não dispõe de configuração de contorno das caches, de modo que este capítulo não reporta ganho de desempenho em ciclos, e sim o custo de cada evento da cache e o efeito da geometria sobre a taxa de acerto. Por sua vez, a síntese lógica foi realizada e é apresentada na seção~\ref{sec:4.8}, mas o projeto não foi executado na placa, de modo que os números de área e de frequência máxima descrevem o que a ferramenta reporta, e não o que se observaria em operação. Registra-se ainda uma diferença de modo entre as duas partes deste capítulo: as medições de ciclo das seções~\ref{sec:4.3.1} a~\ref{sec:4.3.3} usam o arranjo sem blocos de memória dedicados, enquanto a síntese usa o arranjo com eles; a lógica de acerto e de falta é a mesma, mas o custo em ciclos difere, em um ciclo por falta e em dois ciclos por conjunto na drenagem. A quinta é a medição da instrução de sincronização: o custo foi medido, mas sob três restrições que delimitam o alcance dos números das seções~\ref{sec:4.3.1} a~\ref{sec:4.3.3}: a memória do modelo tem latência fixa, sem contenção e sem o comportamento de uma memória dinâmica real, de modo que os valores são em ciclos de relógio e não em tempo; o caso com a cache integralmente suja é um pior caso construído, e não um padrão de uso; e o custo no pipeline foi medido para um programa específico, com variações de alinhamento da instrução seguinte deslocando o termo fixo em poucos ciclos.

\section{Síntese em FPGA}\label{sec:4.8}

\glsadd{FPGA}A verificação estabeleceu que o subsistema faz o que a especificação determina, faltava estabelecer que ele cabe em um dispositivo real e opera na frequência prevista.

A compilação foi feita no Quartus Prime 25.1 Lite, versão 25.1std build 1129, para o dispositivo 5CSEMA5F31C6, o Cyclone V SE da placa DE1-SoC. O projeto sintetizado é o Pipeline3Fpga com o arranjo em blocos de memória dedicados, caches de 8 KiB com quatro vias, linha de 16 bytes e LRU, memória principal de 8192 palavras, ou 32 KiB, com latência um, e restrição de clock de 50 MHz sobre a entrada CLOCK\_50, o que corresponde a um período de 20 ns. A compilação terminou com sucesso, com o Fitter, o Assembler e o Timing Analyzer concluídos.

A \refTab{tab:4.10} reúne o uso de recursos reportado.

{\small
\begin{longtable}{@{}>{\raggedright\arraybackslash}p{(\textwidth - 2\tabcolsep) * \real{0.2383}}>{\raggedright\arraybackslash}p{(\textwidth - 2\tabcolsep) * \real{0.7617}}@{}}
\caption[Uso de recursos na síntese para o Cyclone V da DE1-SoC]{Uso de recursos na síntese para o Cyclone V da DE1-SoC (Fonte: autoria própria, a partir do relatório de compilação de 25 de novembro de 2026)}\label{tab:4.10}\\
\hline
\textbf{Item} & \textbf{Valor} \\
\hline
\endfirsthead
\hline
\textbf{Item} & \textbf{Valor} \\
\hline
\endhead
\hline
\endlastfoot
Lógica (ALMs) & 10.726 de 32.070 (33\%) \\
Registradores & 9.449 \\
Bits de memória em bloco & 414.720 de 4.065.280 (10\%), em 66 blocos M10K \\
Blocos de multiplicação (DSP\glsadd{DSP}) & 9 de 87 (10\%) \\
Pinos & 15 de 457 (3\%) \\
PLLs & 0 de 6 \\
\end{longtable}
}

O resultado de temporização é o que dá sentido ao número de área. A frequência máxima de operação no modelo de pior caso, Slow 1100mV 85C, é de 53,43 MHz, contra os 50 MHz exigidos, o que corresponde a um período de cerca de 18,7 ns e a uma folga de aproximadamente 1,3 ns, já no modelo Slow 1100mV 0C a frequência máxima é de 54,7 MHz, e é o modelo de 85 °C que governa. A margem é estreita, ela indica que o caminho crítico está próximo do limite, de modo que qualquer acréscimo de lógica no caminho que vai do endereço ao sinal de acerto exigiria reavaliação da restrição de clock.

A decomposição do uso por módulo confirma a análise da seção~\ref{sec:2.5.1} e é o resultado mais informativo da síntese. As duas caches somam 7.608 dos 10.726 ALMs, ou 71\% do projeto, e 6.858 dos 9.449 registradores, ou 73\%. Desse total, cada módulo de política de reposição responde por cerca de 1,8 mil ALMs e 1,2 mil registradores, e os arranjos de armazenamento por cerca de 0,9 mil ALMs e 1,4 mil registradores na cache de instruções e 1,9 mil ALMs e 1,9 mil na cache de dados. O restante do processador é pequeno em comparação: a unidade lógica e aritmética responde por 398,7 ALMs, o banco de registradores de controle e status por 252,2 e o banco de registradores por 236,6. Como o armazenamento está nos blocos de memória, esse custo é o da lógica que envolve os arranjos, e é a assinatura quantitativa do argumento da seção~\ref{sec:2.5.1}: em uma cache em vias múltiplas, o que se paga não é o armazenamento, é a comparação e a seleção.

Os bits de memória em bloco fecham a contabilidade e conferem com o que a especificação previa. Os 414.720 bits ocupam 66 blocos M10K e se decompõem em 262.144 bits da memória principal, que são 8192 palavras de 32 bits em 26 blocos, e 152.576 bits das duas caches, em 40 blocos. Cada cache ocupa 76.288 bits, sendo 65.536 de dados e 10.752 de tags, que são 128 conjuntos por 4 vias por 21 bits. A diferença entre esse valor e os 77.312 bits calculados na seção~\ref{sec:3.3} são os 1.024 bits de validade e de sujeira, que residem em registradores e por isso não aparecem como memória em bloco, exatamente como a seção~\ref{sec:3.3} descreve. Por sua vez, a memória principal aparece no relatório como altsyncram, o que confirma que ela foi inferida como memória, e não convertida em registradores.

O histórico da primeira tentativa é a razão de ser de todo o arranjo atual. A compilação anterior, com as memórias de leitura combinacional e uma memória principal menor, de 1024 palavras, estimou 260.684 células lógicas e falhou no Fitter, que exigia 11.441 blocos lógicos contra os 3.207 do dispositivo. A causa está registrada na ferramenta, pois o Quartus não inferiu memória em 40 instâncias, as 20 memórias de cada cache, por reconhecer ali lógica de leitura assíncrona, e o armazenamento foi convertido em 186.880 flip-flops, que responderam por 97\% dos registradores do projeto. Portanto, a leitura combinacional era a diferença entre caber e não caber.

Duas pendências delimitam o alcance desses números. A primeira é que o projeto não foi executado na placa, então a síntese demonstra que ele cabe e fecha temporização, mas não que ele opera. A segunda é a pré-carga do programa na memória principal, que se dá por um arquivo com uma linha de cache por linha, quatro palavras cada, e que só se confirma programando o dispositivo. Assim, os números desta seção são uma garantia de viabilidade, e não uma medida de operação.

Em suma, a síntese lógica fecha o percurso que a seção~\ref{sec:2.1} descreve, em que o subsistema foi especificado, descrito, verificado em simulação e implementado em um dispositivo, e o custo em recursos e a frequência alcançada são os números que faltavam à discussão analítica da seção~\ref{sec:2.5.1}.
